\section{A fixed minimal model of full nef dimension}\label{sec:reduction}

Assume that smooth canonical nonvanishing fails in dimension four.
We first choose a terminal minimal model on which every curve through
a very general point has a fixed positive lower bound for its
canonical degree. We then show that every proper
positive-dimensional subvariety through such a point is of general
type. These two properties will control the moving centers of the
later jet argument. The model remains fixed while its ample
perturbations vary.

\begin{proposition}\label[proposition]{prop:reduction}
If \Cref{thm:nonvanishing} fails, there is a projective
$\Q$-factorial terminal fourfold $Y$, with $K=K_Y$, such that
\begin{enumerate}[label=\textup{(\roman*)}]
\item $K$ is nef, $\kappa(Y,K)=-\infty$, and $K^4=0$;
\item $q(Y)=0$ and the nef dimension of $K$ is four;
\item for some integer $\iota>0$, $\iota K$ is Cartier and
\begin{equation}\label{eq:curve-lower-bound}
 K\cdot C\ge \frac1\iota
\end{equation}
for every integral curve $C$ through a point of a fixed very general
locus in $Y$.
\end{enumerate}
\end{proposition}

\begin{proof}
Start with a smooth counterexample $H$.  Run the canonical minimal
model program.  Existence of the necessary flips follows from
\cite[Corollary~1.4.1]{BCHM2010}.  Termination in the pseudo-effective fourfold case is
\cite[Theorem~1.1]{ChenTsakanikas2023}, applied with zero nef data to
the ordinary pair with zero boundary.  Divisorial contractions lower
the Picard number, so together these statements give termination of
the program.  Pseudo-effectivity is preserved and excludes a Mori
fiber space as its endpoint.  We obtain a projective $\Q$-factorial
terminal minimal model $Y$.

Canonical section spaces in sufficiently divisible degrees are
birationally invariant here.  On a common resolution the differences
from the respective pulled-back canonical divisors are effective
exceptional divisors, and pushing their sheaves forward does not add
sections.  Thus $\kappa(Y,K)=-\infty$.  In particular $K$ is not big;
as it is nef, this says $K^4=0$.

Terminal singularities are rational.  Hence $q(Y)$ agrees with the
irregularity of a resolution.  If it were positive,
\cite[Corollary~4.7]{Fujino2010} would make $K$ semiample, a
contradiction.  Likewise, if the nef dimension were at most three,
\cite[Theorem~4.3]{Ambro2005}, with the established log minimal model
program and log abundance in those dimensions, would make $K$
semiample.  This includes nef dimension zero.  Thus the nef dimension
is four.

The curve property of nef reduction \cite[Theorem~2.1]{BauerEtAl2002} says in this
case that $K\cdot C>0$ for every integral curve through a very general
point.  Choose $\iota$ with $\iota K$ Cartier.  Its degree on an
integral complete curve is an integer, so positivity gives
\Cref{eq:curve-lower-bound}.  This uses only the curve property of nef
reduction, not an Iitaka fibration on $Y$.
\end{proof}

\begin{proposition}\label[proposition]{prop:subvarieties}
For $Y$ as in \Cref{prop:reduction}, every proper positive-dimensional
subvariety through a point of a suitable fixed very general locus is
of general type.
\end{proposition}

\begin{proof}
We first work with one family of subvarieties, and at the end make the
very general locus independent of the family.  Hilbert schemes give
countably many finite-type parameter spaces for integral subvarieties
of $Y$.  Stratification, resolution of the universal families, and
generic smoothness give countably many smooth families
\[
 b:\mathcal V\longrightarrow B,\qquad e:\mathcal V\longrightarrow Y,
\]
whose fibers are smooth projective resolutions of their images and
whose fiber maps are birational onto those images.  To obtain this
description, resolve the total universal family on an integral
parameter stratum, shrink until its fibers are smooth and the fiber
maps birational, and repeat on the complement by noetherian
induction.  Non-dominant evaluation families can be discarded by
excluding their image closures in $Y$.

Consider a dominant family with fibers of dimension $d$, where
$0<d<4$.  We may slice the parameter space to dimension $4-d$ so
that $e$ becomes generically finite and remains dominant.  Indeed, the
parameters incident to a general point of $Y$ have a component of
dimension $\dim B-(4-d)$ on the locus where the fiber maps are
generically embeddings.  General ample cuts in a compactification of
$B$ meet this locus in finitely many points.  We take the cuts very
generally, so a very general point of the slice has the generic value
of every plurigenus of the original family.  There are only countably
many such conditions.

On the resulting smooth total space, ramification and terminality
give
\[
 K_{\mathcal V}\sim_\Q e^*K+E_{\mathcal V},\qquad E_{\mathcal V}\ge0.
\]
For completeness, the rational lift of $e$ to a resolution of $Y$ is
defined at codimension-one points by properness.  The usual
ramification formula there, followed by the effective discrepancies
of the terminal target, proves this divisor inequality.  Its
restriction to a very general fiber is
\begin{equation}\label{eq:fiber-canonical-comparison}
 K_V\sim_\Q e_V^*K+E_V,\qquad E_V\ge0.
\end{equation}
In particular $K_V$ is pseudo-effective.

The dimension of $V$ is at most three.  The minimal model program and
abundance in these dimensions give a good terminal minimal model
$V_{\min}$.  Suppose that $V$ is not of general type.  The semiample
canonical divisor of $V_{\min}$ defines a morphism with connected
positive-dimensional general fibers, on each of which a multiple of
$K_{V_{\min}}$ is trivial.  On a common resolution
$\widetilde V$ of $V$ and $V_{\min}$, there is a complete curve through
a very general point with $K_{\widetilde V}$-degree zero.  Here is the
avoidance needed for this assertion.  The singular locus of the
terminal model and the centers of the exceptional divisors have
codimension at least two.  A general fiber either misses a given
vertical center or meets a dominant center in codimension at least
two.  General sufficiently ample complete intersections in that
fiber through a prescribed very general point give a curve
avoiding all such centers; for a one-dimensional fiber the
fiber itself does so.  Its strict transform has canonical degree
zero by the exceptional discrepancy formula.

Pulling back \Cref{eq:fiber-canonical-comparison} to $\widetilde V$
and adding the effective discrepancies of
$\widetilde V\to V$ gives
\[
 K_{\widetilde V}\sim_\Q \widetilde e^{\,*}K+\widetilde E,
 \qquad \widetilde E\ge0.
\]
Choose the preceding curve through a point outside
$\Supp\widetilde E$ and outside the exceptional locus of the
birational map onto the image subvariety.  It is not contained in
$\widetilde E$, and its image in $Y$ is a curve.  Thus its pulled-back
$K$-degree is at most zero.  Dominance of the total evaluation lets
us choose the point in the very general locus of
\Cref{eq:curve-lower-bound}, which gives a contradiction.

We have proved that the very general fiber of the sliced family is
of general type.  By the choice of slice its plurigenera are the
generic plurigenera of the original smooth family.  Upper
semicontinuity shows that all fibers of that family have at least
these plurigenera, hence are of general type as well.  Finally
exclude the image closures of all non-dominant families.  Countability
of the parameter strata gives the asserted very general locus.
\end{proof}

For the rest of the contradiction argument, fix $Y$, $K$, $\iota$,
and a very general locus on which both propositions apply. We next
choose ample perturbations of $K$ with bounded volume and growing
degree on every curve through this locus. The two propositions will
then turn a section of excessive vanishing order into a contradiction.
